\section{设计报告}

\subsection{需求分析}

本系统面向北斗/GNSS 伪距单点定位实验，设计目标是基于 RINEX 观测文件和广播星历文件，完成从数据导入、卫星位置计算、误差改正、定位解算到结果输出的完整软件流程。系统既要能够通过命令行完成批量解算，也要提供 PyQt 图形界面，便于进行文件选择、参数设置、实时结果查看和轨迹回放。

从实验任务看，系统需要解决的核心问题不是单独完成某一个算法函数，而是把 GNSS 定位中的多个环节串成稳定的数据处理链路。输入端需要兼容普通 RINEX 观测文件和导航文件；算法端需要根据广播星历计算 GPS/BDS 卫星位置、钟差和传播延迟改正；解算端需要根据伪距观测方程迭代求解接收机坐标；输出端需要生成 CSV、误差统计和可视化图像，供后续检查和报告使用。

结合实验文档中的五个核心模块，系统功能需求归纳如下：

\begin{enumerate}
\item 支持 RINEX 观测文件与导航文件解析，提取卫星编号、伪距、信噪比、广播星历参数和测站近似坐标。
\item 基于广播星历手写实现 GPS/BDS 卫星位置计算、卫星钟差修正、相对论效应修正、对流层与电离层延迟改正。
\item 构建伪距定位观测方程，采用迭代最小二乘解算接收机 ECEF 坐标、经纬高坐标和接收机钟差。
\item 支持连续历元定位分析，输出定位轨迹、误差统计、DOP 和卫星可见数变化。
\item 提供命令行与 PyQt GUI 两类入口，实现数据导入、参数设置、实时显示、绘图和轨迹回放。
\end{enumerate}

设计时还需要考虑以下约束。第一，核心定位算法由项目代码手写实现，不调用 RTKLIB、georinex、pymap3d 等现成定位库，第三方依赖主要用于 GUI、绘图和基础运行环境。第二，数据文件可能来自不同 RINEX 版本和不同 GNSS 系统，因此解析模块和解算模块之间需要通过统一数据结构连接。第三，实际数据中可能出现缺少星历、伪距无效、卫星数不足、参数设置错误等情况，系统需要给出可追踪的处理方式，而不是直接中断整个流程。

\subsection{功能设计}

\subsubsection{总体设计思路}

系统采用分层和模块化设计。底层负责 RINEX 文件读取与数据结构封装，中间层负责卫星位置、钟差和传播延迟改正，核心层负责单历元最小二乘定位，流程层负责连续历元循环、结果统计和绘图，最外层提供命令行脚本和 GUI 调用入口。

这种设计方式的主要考虑是让算法模块和界面模块分开。命令行脚本、批量运行脚本和 GUI 都调用同一套解析、修正和定位函数，避免出现“命令行一套逻辑、界面一套逻辑”的情况。后续如果增加新的数据集或新的结果展示方式，也可以尽量复用已有解算流程。

\subsubsection{功能模块划分}

软件按照实验要求划分为五个模块，各模块职责如表 \ref{tab:software-modules} 所示。模块之间的主调用关系为：RINEX 数据解析模块先生成观测历元和星历记录；卫星位置与修正模块根据观测时间和广播星历计算改正后的观测量；单点定位模块完成坐标解算；连续定位模块负责多历元循环和结果保存；系统整合测试模块提供 CLI、GUI 和批量运行入口。

\begin{table}[H]
\centering
\small
\begin{tabular}{p{0.18\linewidth}p{0.24\linewidth}p{0.48\linewidth}}
\toprule
模块 & 主要文件 & 功能说明 \\
\midrule
RINEX 数据解析 & \texttt{rinex\_obs.py}, \texttt{rinex\_nav.py} & 读取 OBS/NAV 文件，解析观测类型、历元数据、伪距、SNR、星历和电离层参数。 \\
卫星位置与修正 & \texttt{satellite.py}, \texttt{atmosphere.py} & 根据广播星历计算卫星 ECEF 坐标、钟差、相对论项、TGD、对流层和电离层改正。 \\
单点定位解算 & \texttt{positioning.py}, \texttt{coords.py} & 构建伪距方程，执行加权迭代最小二乘，输出 ECEF、BLH、PDOP、GDOP 和残差。 \\
连续定位分析 & \texttt{pipeline.py}, \texttt{analysis.py}, \texttt{plotting.py} & 逐历元循环解算，生成 CSV、误差统计、误差曲线、卫星数曲线和轨迹图。 \\
系统整合测试 & \texttt{experiment\_modules.py}, \texttt{scripts/} & 封装五个实验模块，提供 CLI、GUI、批量运行和模型训练入口。 \\
\bottomrule
\end{tabular}
\caption{软件功能模块划分}
\label{tab:software-modules}
\end{table}

\subsubsection{软件架构图}

系统总体架构如图 \ref{fig:software-architecture} 所示。数据从 RINEX 观测文件和导航文件进入后，先经过解析和预处理，再进入卫星位置、钟差和传播延迟改正模块。单点定位模块输出单历元解算结果，连续定位模块进一步生成结果文件和图像，最终由命令行或 GUI 展示。

\begin{figure}[H]
\centering
\resizebox{0.96\linewidth}{!}{%
\begin{tikzpicture}[
  node distance=0.42cm,
  stage/.style={draw, rounded corners, align=center, text width=2.5cm, minimum height=1.05cm, fill=blue!5},
  arrow/.style={-Latex, thick}
]
\node[stage] (a) {RINEX\\OBS/NAV};
\node[stage, right=of a] (b) {数据解析\\与预处理};
\node[stage, right=of b] (c) {卫星位置\\钟差与延迟改正};
\node[stage, right=of c] (d) {SPP\\最小二乘解算};
\node[stage, right=of d] (e) {连续定位\\误差分析};
\node[stage, right=of e] (f) {CSV/图像\\GUI 展示};
\foreach \x/\y in {a/b,b/c,c/d,d/e,e/f}
  \draw[arrow] (\x) -- (\y);
\end{tikzpicture}}
\caption{北斗定位解算软件总体架构}
\label{fig:software-architecture}
\end{figure}

\subsubsection{主要流程图}

定位解算的主要处理流程如图 \ref{fig:positioning-flow} 所示。流程中先完成文件读取和观测预处理，再对每颗可用卫星计算位置和误差改正，最后构建线性化伪距方程并通过迭代最小二乘求解。连续定位时，该流程会按设定步长在多个历元上重复执行。

\begin{figure}[H]
\centering
\begin{tikzpicture}[
  node distance=0.62cm,
  flow/.style={draw, rounded corners, align=center, text width=10.2cm, minimum height=0.62cm, fill=green!4},
  arrow/.style={-Latex, thick}
]
\node[flow] (a) {选择 RINEX 观测文件和导航文件};
\node[flow, below=of a] (b) {解析头文件、历元观测值、广播星历和电离层参数};
\node[flow, below=of b] (c) {剔除无效伪距、不健康卫星、无星历卫星和低高度角卫星};
\node[flow, below=of c] (d) {计算卫星 ECEF 坐标、卫星钟差、地球自转与传播延迟改正};
\node[flow, below=of d] (e) {构建伪距观测方程，按高度角设置权重};
\node[flow, below=of e] (f) {迭代最小二乘求解接收机坐标与钟差};
\node[flow, below=of f] (g) {输出 BLH、PDOP、GDOP、残差和有效卫星数};
\node[flow, below=of g] (h) {连续历元循环，生成 CSV、误差统计和图像};
\foreach \x/\y in {a/b,b/c,c/d,d/e,e/f,f/g,g/h}
  \draw[arrow] (\x) -- (\y);
\end{tikzpicture}
\caption{定位解算全流程}
\label{fig:positioning-flow}
\end{figure}

\subsubsection{数据结构设计}

项目主要数据结构定义在 \texttt{src/models.py} 中。设计时将观测文件、导航文件和定位结果分别封装，避免不同模块直接依赖原始文本文件格式。

\begin{table}[H]
\centering
\small
\begin{tabular}{p{0.28\linewidth}p{0.62\linewidth}}
\toprule
数据结构 & 设计作用 \\
\midrule
\texttt{ObsHeader} & 保存观测文件版本、测站名、近似坐标、观测类型和时间系统。 \\
\texttt{ObsEpoch} & 保存单历元时间、历元标志和以卫星编号为键的观测值字典。 \\
\texttt{NavRecord} & 保存单颗卫星的一组广播星历参数，包括轨道参数、钟差参数、TGD 和健康状态。 \\
\texttt{NavHeader} & 保存导航文件头信息，包括 GPS/BDS 电离层参数和闰秒信息。 \\
\texttt{PositionSolution} & 保存定位结果、接收机钟差、使用卫星数、DOP 和残差诊断指标。 \\
\bottomrule
\end{tabular}
\caption{核心数据结构设计}
\label{tab:data-structures}
\end{table}

这些结构将 RINEX 文件、卫星计算和定位结果解耦，使各模块之间通过明确的数据对象传递信息，便于调试和后续扩展。例如连续定位模块不需要直接解析 RINEX 文本，而是逐历元读取 \texttt{ObsEpoch}；绘图模块也不直接参与定位计算，只读取已经输出的 CSV 结果。

\subsubsection{输入输出与接口设计}

系统对外提供命令行和 GUI 两类接口。命令行入口主要用于复现、批量运行和自动化测试，典型脚本包括 \texttt{scripts/run\_spp.py}、\texttt{scripts/run\_continuous.py}、\texttt{scripts/run\_batch.py} 和 \texttt{scripts/plot\_results.py}。主要输入参数包括观测文件路径、导航文件路径、GNSS 系统类型、最大历元数、历元步长、输出 CSV 路径和图片保存目录。

GUI 入口由 \texttt{scripts/gui\_app.py} 提供，面向交互式演示。界面需要支持观测文件和导航文件选择、系统类型选择、最大历元数和高度角等参数设置，并显示最新定位坐标、使用卫星数、PDOP 和运行状态。运行完成后，GUI 通过绘图和轨迹回放功能查看结果。

输出文件按用途分为三类。第一类是连续定位 CSV，保存每个历元的时间、坐标、误差、DOP、卫星数和残差指标；第二类是汇总表，如 \texttt{results/summary.csv}，用于对不同数据集结果进行比较；第三类是图像文件，包括误差/DOP/卫星数曲线和轨迹散点图。这样的输出设计便于命令行、GUI 和报告图表共用同一批结果。

\subsubsection{异常处理与边界条件设计}

实际 GNSS 数据处理中，经常会遇到观测值缺失、星历不匹配或参数设置不合理的情况。系统在设计上将这些问题放在数据预处理和定位入口处处理，尽量避免错误向后传递。

\begin{enumerate}
\item 当某颗卫星没有有效伪距、星历缺失或健康状态不满足要求时，该卫星不参与当前历元定位。
\item 当卫星高度角低于阈值时，系统剔除该观测，以减小低高度角多路径和大气误差影响。
\item 当可用卫星数不足以构建定位方程时，当前历元不输出有效定位解，但连续定位流程继续处理后续历元。
\item 当步长、最大迭代次数等输入参数不合理时，入口脚本或 GUI 应先给出提示，避免出现除零或无意义循环。
\item 当遇到 Hatanaka 压缩文件等暂不支持格式时，系统识别后提示先转换为普通 RINEX，而不是直接按普通文本解析。
\end{enumerate}

通过这些边界处理，系统可以区分“文件无法解析”“当前历元无法定位”和“定位结果可用但精度较差”等不同情况，便于调试和结果解释。

\subsubsection{可扩展性设计}

当前系统重点支持 GPS、BDS 和 GPS+BDS 联合解算。为了便于后续扩展，系统在设计时尽量避免把系统类型写死在单一函数中，而是通过卫星编号、时间系统和导航记录选择对应处理逻辑。后续如果增加 Galileo、GLONASS 等系统，可以在现有 GNSS 系统识别、导航解析和钟差模型基础上继续扩展。

算法层也保留了继续改进的空间。例如，现有系统使用广播星历和伪距单点定位模型，后续可以进一步引入精密星历、更多电离层/对流层模型、鲁棒估计或动态滤波。界面层与算法层分离后，新增算法功能时不需要重写 GUI，只需要在参数设置和结果展示部分增加相应入口。
